\documentclass[11pt,a4paper]{amsart}
\usepackage[T1]{fontenc}
\usepackage{amssymb}
\usepackage[margin=1in]{geometry}
\usepackage[hidelinks]{hyperref}
\hypersetup{pdftitle={Edge reversal changes Jacobian matroids of homoscedastic tournaments}}
\urlstyle{same}
\DeclareFontShape{OT1}{cmr}{bx}{n}{<-> cmbx10}{}
\DeclareMathSizes{8}{8}{8}{8}
\DeclareMathSizes{9}{9}{8}{8}
\DeclareMathSizes{10.95}{10.95}{8}{8}
\newtheorem{theorem}{Theorem}
\title{Edge reversal changes Jacobian matroids of homoscedastic tournaments}
\date{}
\begin{document}
\begin{abstract}
Wu's dissertation reports computations that imply small counterexamples to
the edge-reversal conjecture for labelled Jacobian matroids. We give explicit
exact certificates and extend the strongly connected counterexamples to
every size \(n\ge4\). On three vertices one edge reversal takes a transitive
tournament to a directed triangle. One precision-coordinate triple has rank
two before the reversal and rank three after it. All twenty triples are
independent for the directed triangle, so these two matroids are also
nonisomorphic as abstract matroids.
\end{abstract}
\maketitle

\section{Question and convention}

For a loopless directed graph \(G\) with no opposite-edge pair, let \(A=I-\Lambda\),
where \(\lambda_{ij}\) is free exactly on arrows \(i\to j\). The precision
parametrization of the homoscedastic linear Gaussian structural equation
model is
\[
 K=\psi_G(s,\Lambda)=sAA^{\top},\qquad s>0,\quad \det A\ne0 .
\]
This is the convention of Drton, Hollering and Wu \cite{dhw}. Its Jacobian
matroid has a ground-set element for each labelled coordinate \(K_{ij}\)
with \(i\le j\), including \(K_{ii}\). A coordinate set has rank equal to
the rank of its derivative matrix over the rational function field.
In the transposed convention of \cite{dhw} the ground-set elements are
columns. In the calculations below coordinates are rows. A nonzero
minor at one admissible point certifies a nonzero generic minor.
The characteristic-zero criterion relating algebraic independence and
Jacobian rank is stated in \cite{dhw}.

A complete simple directed graph in this setting is a tournament.
Conjecture 5.4 of \cite{dhw} asserts equality of the labelled matroids
when two tournaments differ only on the edge between vertices \(n-1,n\).
Simultaneously relabelling both graphs allows us to choose the names
used in the examples before making those endpoints the last two vertices.

\begin{theorem}
For every \(n\ge3\), a single reversal can change the labelled Jacobian
matroid of a homoscedastic tournament. At \(n=3\) the two matroids can also
differ up to abstract isomorphism. For every \(n\ge4\), a labelled
counterexample is possible with both tournaments strongly connected.
\end{theorem}

\section{Three vertices}

Let \(G\) have edges \(2\to1\), \(1\to3\), \(2\to3\). Reverse \(2\to3\)
to obtain the directed cycle \(H\). Write \(a=\lambda_{21}\),
\(b=\lambda_{13}\), and let \(c\) denote the coefficient on the
edge undergoing reversal. The two row matrices are
\[
 A_G=\begin{pmatrix}1&0&-b\\-a&1&-c\\0&0&1\end{pmatrix},
 \qquad
 A_H=\begin{pmatrix}1&0&-b\\-a&1&0\\0&-c&1\end{pmatrix}.
\]
The triple \((u,v,w)=(K_{11},K_{13},K_{33})\) is
\(\bigl(s(1+b^2),-sb,s\bigr)\) for \(G\). For \(H\) the first two
coordinates agree and the third becomes \(s(1+c^2)\).
With \(P=uw-v^2-w^2\), direct substitution gives
\[
 P_G=0,\qquad P_H=s^2c^2(b^2-1-c^2).
\]
For \(G\) the generic rank is at most two, because the triple depends
only on \(s,b\). The derivative matrix for \(v,w\) in \(s,b\) has
determinant \(s\), giving rank exactly two. For \(H\), the determinant
on \((s,b,c)\) is
\[
 \det\begin{pmatrix}
 1+b^2 & 2sb & 0\\[-1pt]
 -b & -s & 0\\[-1pt]
 1+c^2 & 0 & 2sc
 \end{pmatrix}
 = 2s^2 c(b^2-1).
\]
At \(s=1,\ a=b=1/4,\ c=1/2\), the matrix \(A_H\) has determinant
\(1-abc=31/32\). The separator and derivative determinant evaluate
to \(-19/64\) and \(-15/16\), respectively.

\subsection{All triples for the directed cycle}

The six coordinates of \(H\) come in the three pairs
\(\{-sr,s(1+r^2)\}\) for \(r=a,b,c\).
For a triple with one coordinate from each pair, the private-variable
derivatives in \(a,b,c\) give a diagonal minor up to permutation.
Its factors are each of the form \(-s\) or \(2sr\), and their product
is nonzero generically. If a triple takes both coordinates of a pair
with parameter \(r\), the corresponding two-by-two minor on \(s,r\)
is \(s(1-r^2)\) up to sign. The third coordinate has a nonzero
derivative in its own private variable, so the resulting three-by-three
minor is nonzero generically. This treats every triple, since each pair
has only two entries. The number of independent three-element sets
is an invariant of an abstract matroid. \(H\) has all twenty and \(G\)
has at least the one dependent triple exhibited above, proving that
the matroids are nonisomorphic at this size.

For a direct labelled lift to \(n>3\), orient every arrow from a
new vertex to each old vertex, with arbitrary orientations among
new vertices. The old rows have zero entries in the new columns,
so the selected three precision entries and their ranks persist.
Relabel 2 and 3 as \(n-1,n\).
At \(n=2\), the three coordinates in either direction, up to
permutation, are \(s,-sr,s(1+r^2)\). Their three pair minors in
\(s,r\) are nonzero polynomials. Both matroids are therefore uniform
of rank two on three elements, so \(n=3\) is minimal.

\section{Strongly connected counterexamples}

On \(1,2,3,4\), let \(G\) have arrows
\[
 2\to1,\quad3\to1,\quad1\to4,\quad3\to2,\quad4\to2,\quad4\to3 .
\]
Reverse \(3\to2\) to obtain \(H\). The cycles \(1\to4\to3\to2\to1\)
in \(G\) and \(1\to4\to2\to3\to1\) in \(H\) show strong connectivity.
Put \(a=\lambda_{21}\), \(b=\lambda_{14}\), and \(c=\lambda_{23}\)
in \(H\). For \((x,u,y,v)=(K_{11},K_{14},K_{22},K_{12})\), the formula is
\[
 \psi_G:\ (s(1+b^2),-sb,s(1+a^2),-sa),\qquad
 \psi_H:\ (s(1+b^2),-sb,s(1+a^2+c^2),-sa).
\]
Here \(\psi_G,\psi_H\) designate the selected coordinates.
The rank for \(G\) is at most three by its three-parameter formula.
The minor for \((x,u,v)\) on \((s,b,a)\) at \(s=1,a=b=0\)
equals 1. For \(H\), the four-by-four determinant in
\((s,b,a,c)\) equals \(2s^3c(b^2-1)\). At \(s=1,a=b=0,c=1/2\),
with the other coefficients zero, the determinant is \(-1\)
and \(A_H\) is invertible. Both determinant inequalities persist
under sufficiently small perturbations to nonzero edge coefficients.

A second expression makes the dependence visible algebraically.
Let \(D=u^2-v^2\). In \(G\) we have \(s(x-y)=D\) and
\(s^2-xs+u^2=0\), which give
\[
 D^2-xD(x-y)+u^2(x-y)^2=0.
\]
In \(H\), set \(s=1,a=0,b=c=1/2\) and all other coefficients to zero.
The two nonzero edges \(1\to4,2\to3\) give an invertible \(A_H\),
and \(x=y=5/4,\ D=1/4\), so the polynomial equals \(1/16\).

For each new vertex \(t\), orient \(4\to t\) and \(t\to1,2,3\).
Orient the pairs among the new vertices arbitrarily. Each \(t\) can
reach and be reached from the strongly connected core. Rows 1 and 2
keep zeros in every new column. Consequently \(K_{11},K_{22}\)
and \(K_{12}\) stay fixed, as does \(K_{14}\) because each newly
added term has a zero factor in row 1. The coordinates and their
ranks persist. The pair again becomes a counterexample of the
stated conjecture after relabelling 2 and 3 as \(n-1,n\).
This is a labelled rank statement for \(n\ge4\). On three vertices
one reversal of the only strongly connected tournament type destroys
strong connectivity, so the strongly connected pair is minimal.

\section{An earlier three-vertex construction}

Another labelled witness uses two transitive graphs. Keep
\(2\to1,3\to1\) and reverse \(2\to3\). Write
\(a,b,c\) for the corresponding free entries of \(A=I-\Lambda\)
on the varying edge, \(2\to1\), and \(3\to1\).
The triple \((K_{11},K_{13},K_{33})\) is
\((s,sc,s(1+c^2))\) before the reversal and
\((s,sc,s(1+c^2+a^2))\) after it.
The first triple satisfies
\(K_{11}K_{33}-K_{11}^2-K_{13}^2=0\).
The rank is two, since the \((K_{11},K_{13})\) minor on \(s,c\) is \(s\).
The second triple has determinant \(2s^2a\) on \(s,c,a\), so has
rank three. The underlying graphs are isomorphic by exchanging
2 and 3, which induces an isomorphism of their abstract matroids.

\section{Prior computations and the code comparison}

Wu's dissertation \cite[\S5.4.1, p.~63]{wu} reports five labelled
matroids for the eight three-vertex tournaments. The directed triangles
are singleton classes, so a cycle-to-transitive reversal already
contradicts the conjecture by those computations. The reported complete
four-vertex collisions are confined to DAGs, also implying strongly
connected counterexamples at size four. The dissertation was submitted in June and accepted in August 2023.
The first arXiv posting dates to August 2023, and the journal issue to
February 2025.
The direct certificates above prove specific ranks and the separation
of abstract matroids at size three. The strong construction extends
the labelled separation uniformly in size.

The paper reports a Macaulay2 check through six vertices. In the
linked script \cite{code}, line 32 returns only entries indexed by
two-element subsets of the vertex set. Lines 37 and 66--71 use
that list for the matroid comparison. At \(\Lambda=0\) and \(s=1\),
there is exactly one edge variable for each off-diagonal coordinate
of a tournament. The corresponding block of the Jacobian is a
signed permutation matrix. The off-diagonal restriction is thus
free for either orientation. Inclusion of the diagonals is essential
to the conjecture as stated and to the exhibited separators.

\section{Verification and disclosure}

Internally reviewed, not peer reviewed. The exact rational replay
checks the small examples, all twenty triples for each three-vertex
graph, and the strongly connected construction through ten vertices.
The repository contains a partial Lean 4.30.0 check of the selected
three-vertex Gram coordinates, separator identities, polynomial
linearization, minor formula and rational values. The Jacobian
matroid criterion, abstract comparison and uniform extension are
proved in the text. The Lean file does not formalize those results
or the invertibility of the full coefficient matrix.

Under human direction, ChatGPT Astra 6.2 (OpenAI) generated the
discovery proof and its research packet. OpenAI Codex performed the
independent mathematical review and developed the partial Lean
certificate. GPT-6.2 Sol (OpenAI) checked the prior computations
and prepared the manuscript, exact replay and illustrations.
DannyExperiments maintains the repository and is the preferred
citation identity. No AI system is proposed as an author.

\begin{thebibliography}{9}
\bibitem{dhw}
M. Drton, B. Hollering and J. Wu,
\emph{Identifiability of homoscedastic linear structural equation
models using algebraic matroids},
Advances in Applied Mathematics \textbf{163} (2025), Part B, 102794.
\href{https://doi.org/10.1016/j.aam.2024.102794}{doi:10.1016/j.aam.2024.102794}.
\bibitem{wu}
J. Wu, \emph{Homoscedasticity and Feedback Loops in Graphical Models},
doctoral dissertation, Technical University of Munich, accepted 2023.
\href{https://mediatum.ub.tum.de/doc/1713538/document.pdf#page=79}{Section 5.4.1, p.~63}.
\bibitem{code}
M. Drton, B. Hollering and J. Wu,
\emph{Cyclic SEM identifiability}, MathRepo supplementary materials.
\href{https://mathrepo.mis.mpg.de/_downloads/8104bb4d6c771e454c148c4d4e0f542e/CompleteCyclicSEMs.m2}{CompleteCyclicSEMs.m2}.
\end{thebibliography}
\end{document}
